\cotefile{136}
\scanpage{189}\mspage{- 64 -}\hfill IX

une propriété nettement plus forte que celle énoncée dans 5.7, qui pourra alors servir de caractérisation intrinsèque de la partie fixe de $T_p(A_K)$ (en termes de la seule structure de pro-groupe de Barsotti-Tate de $T_p(A_K)$), - du moins dans le cas d'inégales caractéristiques :

\uline{Théorème} 5.8. (M. RAYNAUD). \uline{Soit} $A_K$ \uline{un schéma abélien à réduction semi-stable sur le corps des fonctions} $K$ \uline{d'un trait hensélien} $S$. \uline{Soit} $\handwritten{\ell}$ \uline{un nombre premier} (\uline{le cas intéressant étant celui où} $\handwritten{\ell}=p$), \uline{et considérons le pro-}$\handwritten{\ell}$\uline{-groupe de Barsotti-Tate} $X_{\handwritten{\eta}}=T_{\handwritten{\ell}}(A_K)$ \uline{et sa} ``\uline{partie fixe}'' $X_{\handwritten{\eta}}^f$, \uline{qui est la fibre générique du pro-}$\handwritten{\ell}$\uline{-groupe de Barsotti-Tate} $X^f=T_{\handwritten{\ell}}(A^{\unsure{o|0}})^f$. \uline{Supposons}, \uline{si} $\handwritten{\ell}=p$, \uline{qu'on ait} $\mathrm{car.}\ K=0$ (\uline{de sorte que nous disposons sur} $S$ \uline{du théorème de TATE} [33, th.4]). \uline{Soit} $F$ \uline{un pro-}$\handwritten{\ell}$\uline{-groupe de Barsotti-Tate sur} $S$. \uline{Alors tout homomorphisme} $F_{\handwritten{\eta}}\handwritten{\longrightarrow}X_{\handwritten{\eta}}$ \uline{se factorise par} $X_{\handwritten{\eta}}^f$.

Passant aux ind-$\handwritten{\ell}$-groupes associés aux pro-$\handwritten{\ell}$-groupes de Barsotti-Tate envisagés, l'énoncé précédent devient équivalent au suivant, grâce au théorème cité de TATE :

\uline{Corollaire} 5.9. \uline{Si} $F$ \uline{est un} ind-$\handwritten{\ell}$\uline{-groupe de Barsotti-Tate}, \uline{alors tout homomorphisme} $u_{\handwritten{\eta}}: F_{\handwritten{\eta}}\handwritten{\longrightarrow} A_{\handwritten{\eta}}$ \uline{se prolonge} (\uline{de façon unique}) \uline{en un homomorphisme} $u: F\handwritten{\longrightarrow} A$.

En d'autres termes, la propriété ``néronienne'' caractéristique (1.1.2) du modèle de Néron reste valable pour des homomorphismes de ind-schémas en groupes de Barsotti-Tate, sans hypothèse de lissité sur celui-ci. (Seul le cas des groupes ind-étales résulterait directement de la définition.)

% archive sheet number: 188
% archive numbering: 382 (pencil and printed label, foot).

\scanpage{190}

\[
\boxed{\boxed{\beta(A_1)=T}}\quad\text{i.e.}\quad A_1(\deleted{\unsure{1}}T)=T
\]

\deleted{\unsure{Posons|Donnons}} Considérons

\[
\deleted{\text{\unsure{Don|D'où|Donc}}}\qquad A_1-X_0=\varphi_0^0(X_0)-X_0=-\bigl(\varphi_0^0(X_1)-X_1\bigr)=f\in\mathcal A_1^0
\]

(où $\varphi_0(x)$ est la composante de $\varphi(x)$ suivant

\[
S\{X_0\}\simeq k\widetilde{\{T\}}\{X_0,X_1\}/(X_0+X_1-T)_{pd}\qquad\text{dans}
\]

\[
\mathcal A_1=\underbrace{S\{X_0\}}_{A_1^0}+\underbrace{S\{X_0\}dX_0}_{A_1^1}\simeq S\{X_1\}+S\{X_1\}dX_1\ \bigr)
\]

[ \deleted{NB} Comme $X_1=T-X_0$, on a

\[
\varphi_0^0(X_1)=\varphi_0^0(T)-\varphi_0^0(X_0)\inserted{=T-\varphi_0^0(X_0)},\quad\text{donc}\quad \varphi_0^0(X_1)-X_1=X_0-\varphi_0^0(X_0)\ \bigr]
\]

Alors la conjonction des conditions $A_1(X_0)=0$ et $A_1(T)=T$ s'écrit aussi

\[
 f\in(X_0)_{pd}\cap(X_1)_{pd}
\]

donc, si \unsure{l'on|on} interprète $\mathcal A_1^0$ comme un \deleted{\unsure{quotient de|produit de}} $k\{X_0,X_1\}$, cela signifie qu'on doit avoir

\[
 f=\sum_{i,j\geq\underline{\underline{1}}}\deleted{\unsure{a}}f_{ij}X_0^{(i)}X_1^{(j)}
 \qquad f_{ij}\in k
\]

d'où en résumé \deleted{$A_1=$}

\uline{Prop}. Les endom. $\varphi$ de la $S$-algèbre diff\textsuperscript{lle} $\mathcal A_1\simeq S\{X_0,X_1,dX_0,dX_1\}/(X_0+X_1-T,dX_0-dX_1)_{pd}$, (où $S=k\{T\}$ avec ses puiss. div) qui sont compatibles aux deux augm. $\alpha,\beta$ $\mathcal A_1\unsure{-|\to} S$ ($\alpha(X_0)=0,\alpha(X_1)=T$; $\beta(X_0)=T,\beta(X_1)=0$) sont donnés exactement par les systèmes

\[
 f\in\mathcal A_1^0=k\{X_0,X_1\}
\]

avec $f\in\operatorname{Ker}(k\{X_0,X_1\}\unsure{-|\to}k\{X_0\})\cap\operatorname{Ker}(k\{X_0,X_1\}\unsure{-|\leftarrow|\to}k\{X_1\})$

i.e. $\boxed{f=\sum_{i,j\geq1}a_{ij}X_0^{(i)}X_1^{(j)}}$

% archive sheet number: 189

\scanpage{191}\mspage{- 62 -}\hfill IX

où on se rappellera que les groupes $T_{\handwritten{\ell}}(A_K)^f$ et $T_{\handwritten{\ell}}^{\unsure{1|\prime}}(A'_K)^f$ qui figurent au second membre sont les restrictions à $\handwritten{\eta}$ de pro-$\handwritten{\ell}$-groupes de Barsotti-Tate définis sur $S$, $T_{\handwritten{\ell}}(A^{\unsure{o|0}})^f$ et $T_{\handwritten{\ell}}(A^{\prime o})^f$ respectivement ; la lettre $D$ désigne ici les duals au sens des pro-groupes de Barsotti-Tate, savoir $\underline{\mathrm{Hom}}(-,\mathbb{Z}_{\handwritten{\ell}}(1))$. Ainsi :

\uline{Proposition} 5.6. \uline{Supposons} $S$ \uline{complet et que} $A_K$ \uline{ait réduction semi-stable sur} $S$. \uline{Alors pour tout nombre premier} $\handwritten{\ell}$ (\uline{qui peut être égal à la caractéristique résiduelle de} $S$), \uline{le pro-}$\handwritten{\ell}$\uline{-groupe de Barsotti-Tate} $T_{\handwritten{\ell}}(A_K)^t$ \uline{est la fibre générique d'un pro-groupe de Barsotti-Tate sur} $S$, ``\uline{de type torique}'', \uline{donné par} (5.5.7), \uline{tandis que} $T_{\handwritten{\ell}}(A_K)/T_{\handwritten{\ell}}(A_K)^f$ \uline{est la fibre générique d'un pro-groupe de Barsotti-Tate sur} $S$ \uline{de type pro-étale}, \uline{donné par} (5.5.8). \uline{Enfin les groupes} $T_{\handwritten{\ell}}(A_K)^f$ \uline{et} $T_{\handwritten{\ell}}(A_K)/T_{\handwritten{\ell}}(A_K)^t$ \uline{sont les fibres génériques de pro-groupes de Barsotti-Tate définis sur} $S$, \uline{égaux respectivement à} $T_{\handwritten{\ell}}(A^o)^f$ \uline{et à} $D(T_{\handwritten{\ell}}(A^{\prime o})^f)=\underline{\mathrm{Hom}}\,(T_{\handwritten{\ell}}(A^{\prime o})^f,\mathbb{Z}_{\handwritten{\ell}}(1))$.

5.6.1. On notera qu'il revient au même pour le revêtement étale $S'$ de déployer $T_{\handwritten{\ell}}(A_K)^t$ ou de déployer $T_{\handwritten{\ell}}(A_K)/T_{\handwritten{\ell}}(A_K)^f$, car cela signifie respectivement que $S'$ déploye les groupes constants tordus $\underline{M}$ et $\underline{M}'$, or ceux-ci sont isogènes, puisque $T_o$ et $T'_o$ le sont. On voit en même temps qu'il existe une plus petite extension étale $S'$ de $S$ qui fait l'affaire.

Signalons aussi la conséquence facile suivante des définitions de 2.2.3 :

% archive sheet number: 190
% archive numbering: 380 (pencil and printed label, foot).

\scanpage{192}

et les

\[
\boxed{B_i\in\mathcal A_1^0=k\{X_0,X_1\}\qquad (i\geq0)}
\]

par les formules

\[
\left\{\begin{aligned}
\varphi(X_0^{(i)})&=(X_0+f)^{(i)}+B_i\,dX_0\\
\varphi(dX_0)&=1+df
\end{aligned}\right.
\]

(de sorte qu'on a

\[
\begin{aligned}
 f&=\varphi_{\deleted{\unsure{0}}}^{0}(X_0)-X_0=-\bigl(\varphi_{\deleted{\unsure{0}}}^{0}(X_1)-X_1\bigr)\\
 B_i\,d\deleted{\unsure{f}}X_0^{\unsure{\prime}}&=\varphi^1\bigl(X_0^{(i)}\bigr)
\end{aligned}
\]

\uline{NB}. Si on échange le rôle de $X_0$ et $X_1$, $f$ change de signe, - que deviennent les $B_i$ ? Il faut calculer

\[
\begin{aligned}
\varphi^1(X_{\deleted{\unsure{0}}1}^{(i)})&=\varphi^1\bigl((T-X_0)^{(i)}\bigr)\\
&=\varphi^1\biggl(\sum_{0\leq j\leq i}T^{(i-j)}X_0^{(j)}(-1)^j\biggr)
=\biggl(\sum_{0\leq j\leq i} T^{(i-j)}(-1)^j B_j\biggr)dX_0\\
&=\underbrace{\deleted{\unsure{-}}\biggl(\sum_{0\leq j\leq i}(-1)^{j+1} T^{(i-j)} B_j\biggr)}_{C_i}dX_1
\end{aligned}
\]

\deleted{\unsure{En termes|En tout cas}} \deleted{$\varphi^1(T)$}

\[
\begin{aligned}
&\underbrace{\varphi^1}_{\unsure{\pi_0^1\varphi|k'_0\varphi}}
\biggl(\underbrace{(X_0+X_1^{\unsure{\prime}})^{(i)}}_{\displaystyle\sum_{j+k=i}X_0^{(j)}X_1^{(k)}}\deleted{\unsure{=}}\biggr)
=\varphi^1(T^i)=0\\
&\underbrace{\hspace{4.4cm}}_{\displaystyle\pi^1\sum_{j+k=i}
\underbrace{\varphi(X_0^{(j)})}_{(X_0+f)^{(j)}+B_jdX_0}
\underbrace{\varphi(X_1^{(k)})}_{(X_1-f)^{(k)}+C_k\underbrace{dX_1}_{-dX_0}}}
\end{aligned}
\]

donc

\[
\sum_{j+k=i}\bigl[-(X_0+f)^{(j)}C_k\deleted{\unsure{-}}+B_j(X_1-f)^{(k)}\bigr]=0
\]

% archive sheet number: 191

\scanpage{193}\mspage{- 61 -}\hfill IX

dans ce cas

\[
M=\mathrm{Hom}_{k\text{-}\mathrm{gr}}(T'_{\unsure{o|0}},\mathbb{G}_{mk})\ ,\tag*{(5.5.4)}
\]

et on aura alors

\[
\underline{M}=M_S\ .
\]

Ceci posé, on aura évidemment, pour tout entier $n$, des isomorphismes

\[
\underline{M}'_n\simeq D(T_n)\ ,\qquad \underline{M}_n\simeq D(T'_n)\ ,\tag*{(5.5.5)}
\]

où $T'_n$ est défini en termes de $A'$ comme $T_n$ en termes de $A$ (5.1.2), d'où on déduit des isomorphismes canoniques (où on pose $\mathbb{Z}_{\handwritten{\ell}}(1)=T_{\handwritten{\ell}}(\mathbb{G}_{mS})$) :

\[
T_{\handwritten{\ell}}(T_n)\simeq\check{\underline{M}}'_n\otimes_{\mathbb{Z}}\mathbb{Z}_{\handwritten{\ell}}(1)\ ,\qquad
T_{\handwritten{\ell}}(T'_n)\simeq\check{\underline{M}}_n\otimes_{\mathbb{Z}}\mathbb{Z}_{\handwritten{\ell}}(1)\ ,\tag*{(5.5.6)}
\]

et par passage à la limite sur $n$ des isomorphismes canoniques :

\[
\handwritten{\left\{}\begin{aligned}
T_{\handwritten{\ell}}(A^o)^t&=T_{\handwritten{\ell}}(\hat{T})\simeq\check{\underline{M}}'\otimes\mathbb{Z}_{\handwritten{\ell}}(1)\\
T_{\handwritten{\ell}}(A^{\prime o})^t&=T_{\handwritten{\ell}}(\hat{T}')\simeq\check{\underline{M}}\otimes\mathbb{Z}_{\handwritten{\ell}}(1)\ .
\end{aligned}\right.\tag*{(5.5.7)}
\]

Dans le cas de la réduction semi-stable, on en conclut par dualité, en utilisant le théorème d'orthogonalité 5.2 et la relation (5.2.2), des isomorphismes canoniques :

\[
\handwritten{\left\{}\begin{aligned}
T_{\handwritten{\ell}}(A_K)/T_{\handwritten{\ell}}(A_K)^{\deleted{\unsure{t}}f}&\simeq\underline{M}_K\otimes_{\mathbb{Z}}\mathbb{Z}_{\handwritten{\ell}}\\
T_{\handwritten{\ell}}(A')/T_{\handwritten{\ell}}(A'_K)^f&\simeq\underline{M}'_K\otimes_{\mathbb{Z}}\mathbb{Z}_{\handwritten{\ell}}\ .
\end{aligned}\right.\tag*{(5.5.8)}
\]

On trouve de même, grâce à la même formule (5.2.2), des isomorphismes canoniques :

\[
\handwritten{\left\{}\begin{aligned}
T_{\handwritten{\ell}}(A_K)/T_{\handwritten{\ell}}(A_K)^t&\simeq D(T_{\handwritten{\ell}}(A'_K)^f)\\
T_{\handwritten{\ell}}(A'_K)/T_{\handwritten{\ell}}(A'_K)^t&\simeq D(T_{\handwritten{\ell}}(A_K)^f)\ ,
\end{aligned}\right.\tag*{(5.5.9)}
\]

% archive sheet number: 192
% archive numbering: 379 (pencil and printed label, foot).

\scanpage{194}

\deleted{\unsure{1}} \uline{Prop} Soit $\varphi$ un endom de $\deleted{A_1=k\{X_0,X_1,dX_0,dX_1\}/d\unsure{X_0+}}$ $\mathcal A_1=k\{X_0,X_1\}[dX_0,dX_1]/(dX_0\unsure{+|-}dX_1)_{\deleted{\unsure{pd}}}$, et $\deleted{\unsure{f_1}}\inserted{\varphi^{\unsure{\prime}}}$ un endom de $S$-algèbre différentielle ($S=k\{T\}$) compatible aux deux augmentations $\alpha$ et $\beta$ (qui en fait sont les deux applications simpliciales $\mathcal A_1\to\mathcal A_0$ !). Soit $n\geq0$ ; je dis qu'il existe un et un seul endom de $S$-Alg. différentielle de $\mathcal A_n$, telle que pour les applications simpliciales (\deleted{\unsure{surj}} croissantes !) $\sigma:\Delta_n\unsure{-|\to}\Delta_1$, l'hom. $\mathcal A_*(\sigma):\mathcal A_1\unsure{-|\to}\mathcal A_n$ satisfasse

\[
\deleted{\unsure{\varphi_*\mathcal A_*}}\qquad\varphi_n\mathcal A_*(\sigma)=\mathcal A_*(\sigma)\varphi_1\ .
\]

\deleted{Si $n$} (Si $n\geq1$, on peut se \deleted{\unsure{supposer}} \inserted{borner} dans cette considération aux $\sigma$ croissantes).

% margin note written sideways
\margin{faux tel quel ; vrai ssi $2B_i=0\quad\forall i\geq1$ i.e. $2\varphi^1(X_0^{(i)})=0\quad\forall i\geq1$ (\deleted{\unsure{bien}} \unsure{vide si|évidt si|Ceci est vrai si} $2\,1_k=0$ ; cela signifie, \deleted{\unsure{que|pour}} \deleted{$\unsure{B_i}$} $2.1_k$ \unsure{rég|inv.}$^{\mathrm r}$) que $B_i=0\quad\forall i\geq1$ i.e. $\varphi$ conserve le degré (et alors $\varphi\unsure{\Leftrightarrow|\Rightarrow}$ un \deleted{\unsure{endom}} de $\mathcal A_1^0$ qui soit comp. aux $\alpha^0,\beta^0$).}

\uline{Cor} \deleted{On a $\varphi_1=\varphi$, et} le système $\varphi_*=(\varphi_n)$ est un endom de la $S$-algèbre différentielle semi-simpliciale $\mathcal A_*$, \uline{caractérisé} par le fait que $\varphi_1=\varphi$. L'ens des endom de $S$-alg. diff $\mathcal A_*$ est ainsi en correspondance bi-univoque (conservant les compositions) avec l'ens des $\varphi$ \inserted{qui satisfont aux conditions corrigées plus haut}
% slanted
(à partir de la prop).

\uline{Dém. du corollaire} Il suffit de montrer que si $\tau:\Delta_{n'}\unsure{-|\to}\Delta_n$, d'où $\mathcal A_*(\tau):\mathcal A_n\unsure{-|\to}\mathcal A_{n'}$

\[
\varphi_{n'}\mathcal A_*(\tau)=\mathcal A_*(\tau)\varphi_n\ .
\]

\inserted{C'est clair si $n=0$ car $\mathcal A_0=S$, et par hyp. si $n=1$.} Si on a $n\geq1$ \unsure{et si on|ou bien} \inserted{alors} regarde les $\sigma:\Delta_n\to\Delta_1$ \deleted{\unsure{qui seront}} (on peut se borner aux surjectifs) \inserted{on verra \unsure{que}} les $\mathcal A_*(\sigma):\mathcal A_1\unsure{-|\to}\mathcal A_n$ \deleted{\unsure{ont|sont}} des images qui \uline{engendrent} l'algèbre $\mathcal A_n$, donc il suffit de prouver que

% archive sheet number: 193

\scanpage{195}\mspage{- 60 -}\hfill IX

\uline{i}.\uline{e}. \uline{définit un isomorphisme de chacun de ces schémas abéliens avec le dual de l'autre}.

La démonstration de 5.4 dans le cas général va être donnée au n$^\circ$ 7, comme conséquence d'une construction fort naturelle et importante due à RAYNAUD. Notons seulement ici que pour prouver 5.4, on peut supposer $S$ complet, compte tenu de 3.3.

5.5. Pour exprimer commodément les facteurs extrêmes dans les filtrations (5.1.10) et (5.1.11), dans le cas de la réduction semi-stable, introduisons les schémas en groupes constants tordus sur $S$

\[
\underline{M}'_{\unsure{o|0}}=D(T_o)\ ,\qquad\underline{M}_o=D(T'_o)\ ,\tag*{(5.5.1)}
\]

duals des tores maximaux $T_o$, $T'_o$ des fibres spéciales de $A,A'$ (SGA 3 X 5.7). Comme $S$ est local complet, ces schémas en groupes peuvent être considérés comme les fibres spéciales de schémas constants tordus

\[
\underline{M}'\ ,\ \underline{M}\text{ sur }S\ ;\tag*{(5.5.2)}
\]

ainsi, si le tore $T_o$ est déployé (par exemple si $k$ est séparablement clos) i.e. si $\underline{M}'_o$ est le groupe constant défini par le groupe ordinaire

\[
M'=\mathrm{Hom}_{k\text{-}\mathrm{gr}}(T_o,\mathbb{G}_{mk})\ ,\tag*{(5.5.3)}
\]

alors $\underline{M}'$ est le groupe constant défini par $M'$ :

\[
\underline{M}'=M'_S\ .
\]

On peut d'ailleurs noter que, puisque $T_o$ et $T'_o$ sont isogènes (2.2.6), $T'_o$ est déployé si $T_o$ l'est, de sorte qu'il y a lieu d'introduire également

% archive sheet number: 194
% archive numbering: 378 (pencil and printed label, foot).

\scanpage{196}

\[
\begin{aligned}
\underbrace{\varphi_{n'}\underbrace{\mathcal A_*(\tau)\mathcal A_*(\sigma)}_{\mathcal A_*(\sigma\tau)}}_{\mathcal A_*(\sigma\tau)\varphi\;(\text{p. hyp. sur }\varphi_{n'})}
&=\underbrace{\mathcal A_*(\tau)\underbrace{\varphi_n\mathcal A_*(\sigma)}_{\mathcal A_*(\sigma)\varphi\;(\text{par hyp. sur }\varphi_n)}}_{\mathcal A_*(\sigma\tau)\varphi}
\qquad\forall\sigma:\Delta_n\to\Delta_1
\end{aligned}
\]

OK.

\uline{Dém. de la prop} Le cas $n=0$ est trivial, pour le cas $n=1$ il faut considérer les trois cas $\sigma=\mathrm{id}$ qui donne $\varphi_1=\varphi$, puis $\sigma$ application constante de valeur $0$ ($\alpha_0$) ou $1$ ($\deleted{\unsure{\alpha}}\beta_0$), de sorte que $\mathcal A_*(\alpha_0)=\alpha$, $\mathcal A_*(\beta_0)=\beta$, et le fait que $\varphi_1=\varphi$ y soit comp. ce n'est autre que l'hyp. suppl. sur $\varphi$ sur les augmentations \inserted{$\alpha,\beta$}. On peut donc supposer $n\geq2$. Considérons les applications \inserted{\unsure{croiss|corresp.}} $\sigma:\Delta_n\to\Delta_1$ ; elles sont surj. ou constantes. Considérons les surjections, il y en a $n$, soit $\sigma_{\deleted{\unsure{\alpha}}r}$ ($0\leq\deleted{\unsure{\alpha}}r\leq n-1$)

\[
\sigma_{\deleted{\unsure{\alpha}}r}(0)\overset{\inserted{\sigma}}{=}0,\,-\,\sigma_{\deleted{\unsure{\alpha}}r}(\unsure{\alpha|r})=0,\quad
\sigma_{\deleted{\unsure{\alpha}}r}(\unsure{\alpha|r}+1)=1,\,-\,\sigma_r(n)=1
\]

donc \deleted{$\mathcal A$}

\[
\left\{\begin{aligned}
\mathcal A_*(\sigma_r)(X_0)&=\deleted{\unsure{U}}X_0+-+X_r\overset{\text{déf}}{=}U_r\\
\mathcal A_*(\sigma_r)(X_1)&=X_{r+1}\,-+X_n
\end{aligned}\right.
\]

On doit donc avoir

\[
\mathcal A_*(\sigma_r)(X_0^{(i)})=U_r^{(i)}
\]

\uline{Lemme} Les $U_r^{(i)}$ ($\deleted{\unsure{0\leq}}i\geq1,\ 0\leq r\leq n-1$) engendrent la \inserted{$S$-}algèbre $\mathcal A_n^0$, qui est isomorphe en fait à $\deleted{\unsure{k}}S\{U_0,-,U_{n-1}\}$,

% archive sheet number: 195

\scanpage{197}\mspage{- 59 -}\hfill IX

désignant respectivement les parties abéliennes de $A_{\unsure{o|0}}^o$ et $A^{\prime o}_o$. Comme, en vertu de VIII 4.8, la biextension $W_o$ provient d'une biextension $W_o^{\mathrm{ab}}$ de $(B_o,B'_o)$ par $\mathbb{G}_{mk}$, déterminée à isomorphisme unique près, l'accouplement (5.3.5) peut aussi s'interpréter comme déduit de la correspondance divisorielle $W_o^{\mathrm{ab}}$, et l'homomorphisme correspondant

\[
T_{\handwritten{\ell}}(B_o)\handwritten{\longrightarrow} D(T_{\handwritten{\ell}}(B'_o))\simeq T_{\handwritten{\ell}}(D(B'_o))\tag*{(5.3.6)}
\]

(où dans le dernier terme, la lettre $D$ désigne le schéma abélien dual) n'est autre que celui déduit par fonctorialité de $T_{\handwritten{\ell}}$ de l'homomorphisme

\[
B_o \handwritten{\longrightarrow} D(\unsure{\mathbb{B}|B}'_o)\tag*{(5.3.7)}
\]

défini par $W_o^{\mathrm{ab}}$. Donc le fait que (5.3.5) soit une dualité parfaite signifie aussi, par le dictionnaire bien connu (cf. [ 5]) que le noyau de (5.3.7), qui est un groupe fini sur $k$ en vertu de 1.5, a une composante $\handwritten{\ell}$-primaire nulle. Dire qu'il en est ainsi pour tout $\handwritten{\ell}$ signifie donc que l'isogénie (5.3.7) est un isomorphisme, i.e. que la correspondance divisorielle $W_o^{\mathrm{ab}}$ est une dualité parfaite entre $B_o$ et $B'_o$. Les considérations qui précédaient montrent seulement que le noyau de (5.3.7) est un $p$-groupe. Enonçons donc le résultat plus complet :

\uline{Théorème} 5.4. \uline{Soit} $S$ \uline{un trait}, \uline{et} $A_K$ \uline{un schéma abélien sur} $S$ \uline{ayant une réduction semi-stable sur} $S$. \uline{Désignons par} $B_o$, $B'_o$ \uline{les parties abéliennes des modèles de Néron de} $A_K$ \uline{et du schéma abélien dual} $A'_K$ \uline{respectivement}. \uline{Alors la correspondance divisorielle canonique} $W_o^{\mathrm{ab}}$ \uline{sur} $(B_o,B'_o)$, \uline{déduite par} VIII 4.8 \uline{de la biextension canonique} $W_o$ \uline{de} $(A_o^o,A^{\prime o}_o)$ (\uline{fibre spéciale de la biextension} $W$ \uline{de} 1.4), \uline{est une dualité parfaite entre} $B_o$ \uline{et} $B'_o$,

% archive sheet number: 196
% archive numbering: 377 (pencil and printed label, foot).

\scanpage{198}

\deleted{\unsure{donc}} (donc c'est la \inserted{$S$-}algèbre engendrée par les $U_{r,i}=U_r^{(i)}\quad 0\leq r\leq n-1,\ i\geq1$, soumis aux relations $U_{r,i}U_{r,j}=C_{ij}U_{r,i+j}$). \unsure{De $=$|Et} on a $\mathcal A_n\simeq S\{U_0,-,U_{n-1}\}[dU_0,-,dU_{n-1}]$

C'est immédiat, car $X_0,-,X_{\unsure{r|n}-1}$ se réconstituent par $U_0,-,U_{n-1}$ par

\[
X_0=U_0,\ X_1=U_1-U_0,\ --,\ X_{n-1}=U_{n-1}-U_{n-2}
\]

et on a $\mathcal A_n^0\simeq S\{X_0,-,X_{n-1}\}$. \unsure{Or}

Ceci montre l'unicité de $\varphi_n$, et qu'on devra avoir

\[
\begin{aligned}
\varphi_n(\underbrace{U_{r,i}^{\deleted{\unsure{(i)}}}}_{U_r^{(i)}})
 &=\varphi_n\bigl(\underbrace{\sigma_r^*}_{\substack{\Vert\,\text{déf}\\\mathcal A_*(\sigma_r)}}(X_0^{(i)})\bigr)
 =\sigma_r^*\bigl(\varphi(X_0^{(i)})\bigr).\\
 &=\sigma_r^*(\xi_i)\ :
\end{aligned}
\]

i.e.

\[
\varphi_n(U_{r\deleted{\unsure{,i}}}^{(i)})=\sigma_r^*(\xi_i)
=\biggl(\deleted{\unsure{A_r U}}\sum_j a_{ij}U_r^{(j)}+\Bigl(\sum b_{ij}U_r^{(j)}\Bigr)dU_r\biggr)
\]

Il y a \inserted{donc} un hom. d'algèbres \inserted{unique} $\mathcal A_n^0\overset{\varphi_n^0}{\unsure{-|\to}}\mathcal A_n$ tel que

\[
\varphi_n^0(U_r^{(i)})=\sigma_r^*(\xi_i)\overset{\text{déf}}{=}\xi_{r,i}
\]

car cela revient à vérifier les relations

\[
\deleted{\sigma_r^*(\xi_i)\sigma_r^*(}\qquad
\xi_{r,i}\xi_{r,j}=C_{ij}\xi_{r,i+j}
\]

i.e.

\[
\begin{aligned}
\underset{\substack{\Vert\\\sigma_r^*(\xi_i\xi_j)}}{\sigma_r^*(\xi_i)\sigma_r^*(\xi_j)}&=C_{ij}\sigma_r^*(\xi_{i+j})\\
\sigma_r^*(C_{ij}\xi_{i+j})&=C_{ij}\sigma_r^*(\xi_{i+j})\qquad\mathrm{OK}
\end{aligned}
\]

% archive sheet number: 197

\scanpage{199}\mspage{- 58 -}\hfill IX

qui est une \uline{dualité parfaite}, i.e. identifiant chacun des deux facteurs du premier membre au dual de l'autre (au sens de 5.2.2). (NB les exposants ab désignent les ``parties abéliennes'', quotient de la partie fixe par la partie torique). Il est clair d'autre part que (5.3.2) est induit par l'accouplement

\[
T_{\handwritten{\ell}}(A^{\unsure{o|0}})^{\mathrm{ab}}\times T_{\handwritten{\ell}}(A^{\prime o})^{\mathrm{ab}}\handwritten{\longrightarrow} T_{\handwritten{\ell}}(\mathbb{G}_{mS})\tag*{(5.3.3)}
\]

déduit par passages aux quotients de l'accouplement (5.2.4.2). Celui-ci étant une dualité parfaite sur les fibres génériques, i.e. l'homomorphisme correspondant

\[
T_{\handwritten{\ell}}(A^o)^{\mathrm{ab}}\handwritten{\longrightarrow} D(T_{\handwritten{\ell}}(A^{\prime o})^{\mathrm{ab}})\tag*{(5.3.4)}
\]

étant un isomorphisme sur les fibres génériques, nous avons évidemment envie d'en déduire qu'il en est de même sur $S$ tout entier, ou ce qui revient au même, sur les fibres spéciales. Cette déduction est valable, trivialement, si $\handwritten{\ell}\handwritten{\ne}p$ ; elle est valable également, du moins dans le cas d'inégales caractéristiques, si $\handwritten{\ell}=p$, grâce à un théorème profond de TATE [33, th.4]. Nous allons cependant pouvoir l'établir directement, sans l'aide du théorème de Tate, plus bas (7.4.6).

Pour préciser la signification géométrique du résultat que nous désirons établir, notons que l'accouplement déduit de (5.3.3) par passage aux fibres spéciales n'est autre que l'accouplement

\[
T_{\handwritten{\ell}}(B_o)\times T_{\handwritten{\ell}}(B'_o)\handwritten{\longrightarrow} T_{\handwritten{\ell}}(\mathbb{G}_{mk})\tag*{(5.3.5)}
\]

déduit de la biextension $W_o$ de $(A_o^o,A^{\prime o}_o)$ par passage au quotient, $B_o$ et $B'_o$

% archive sheet number: 198
% archive numbering: 376 (pencil and printed label, foot).

\scanpage{200}

Si on veut étendre $\varphi_n^0$ en \deleted{$\varphi_n$} : hom de $S$-algèbre $\varphi_n:\mathcal A_n\unsure{-|\to}\mathcal A_n$, il faut se donner les images des $dU_0,-,dU_{n-1}$, et vérifier seulement que ce sont des éléments de \uline{carré nul}, \deleted{\unsure{de degré impair}} \inserted{qui anticomm$^{\mathrm t}$ \deleted{\unsure{deux à deux}} entre eux et commutent} aux $\xi_{s,i}$

\[
\varphi_n(dU_r)=d\varphi_n(U_r)=d\sigma_r^*(\xi_1)=\sigma_r^*(d\xi_1)
\]

\[
\begin{aligned}
 d\xi_1&=\sum_{j\geq1}a_{1j}X_0^{(j-1)}dX_0
 &\biggl(&=d\Bigl(\sum_{j\geq0}a_{1j}U_r^{(j)}\Bigr)\\
 \deleted{d\xi_1}&\deleted{=\sum_{j\geq1}a_{1j}X_0^{(j-1)}dX_0}
 &&=\Bigl(\sum_{j\geq1}a_{1j}U_r^{(j-1)}\Bigr)dU_r\biggr)
\end{aligned}
\]

Ces éléments sont bien de carré nul, anticommutent entre eux, mais commutent-ils aux $\xi_{s,i}=\sigma_s^*(\deleted{\unsure{\xi}}X_0^{(i)})=\sum_{j\geq0}a_{ij}U_s^{(j)}+(\sum_{j\geq0}b_{ij}U_s^{(j)})dU_s$

\deleted{\unsure{En résumé, écrivons|En résumé, écrivant} $\mathcal A_n=S\{U_0,-,U_{n-1}\}[dU_0,-,dU_{n-1}]$ et pour l'élém\textsuperscript{t} général $=F+\sum_{0\leq r\leq n-1}G_r\,dU_r$.}

% slanted
\margin{et $\sigma_r^*(d\xi_1)$ commute bien aux $\xi_{si}=\sigma_s^*(X_0^{(i)})$ pour $s=r$, mais pour $s\ne r$ c'est vrai ssi $2\bigl(\sum_{j\geq1}a_{1j}U_r^{(j-1)}\bigr)\bigl(\sum_{j\geq0}b_{ij}U_s^{(j)}\bigr)=0$, i.e. $\boxed{\substack{2a_{1,j}b_{i\deleted{j}k}=0\\\forall j\geq1,\ k\geq0}}$ Or $\unsure{\alpha}a_{11}=1+\sum_{i\geq\deleted{\unsure{0}}1}\alpha_i T^{(i)}$ est \unsure{non|un} div de \unsure{Zéro|Zéros}, donc il faut $2b_{ik}=0\quad(\forall i\geq1,k\geq0)$ donc $2B_i=0$.}

On utilise, pour prouver que $\varphi_n$ (qui à priori est un hom de $S$-algèbres) commute à la différentielle, on utilise le

\uline{Lemme} $S$ anneau commutatif, $U_0-U_m$ des variables, on considère la $S$-algèbre diff $\mathcal A\deleted{\unsure{\text{anneau}}}=S\{U_0,-,U_m\}[dU_0,-,dU_m]$ (\uline{NB} on néglige la structure des p.d, \unsure{au moins|du moins} $dU_r^{(i)}=U_r^{(i-1)}dU_r$ --). Soit $B$ une $S$-alg. diff. \unsure{alternée}, (grad. mod 2, différentielle de degré $+1$), On s'intéresse aux hom de \inserted{$S$-}algèbres diff $\mathcal A\xrightarrow{\varphi} B$. Pour les \unsure{dét.|tels|tel}, on a

% archive sheet number: 199

\scanpage{201}\mspage{- 57 -}\hfill IX

5.2.6. Supposons maintenant que $A_K$ ait une réduction semi-stable sur $S$, de sorte que nous pouvons prendre pour $\handwritten{\ell}$ n'importe quel nombre premier (distinct ou non de $p$). D'ailleurs, en vertu du critère galoisien de réduction semi-stable 3.5, on a pour $\handwritten{\ell}\handwritten{\ne}p$

\[
(T_{\handwritten{\ell}}(A'_K)^f)^{\handwritten{\perp}}\quad(=T_{\handwritten{\ell}}(A_K)^t)\ \handwritten{\subset}\ T_{\handwritten{\ell}}(A_K)^f\ .\tag*{(5.2.7)}
\]

Il reste à prouver, pour terminer la preuve de 5.2, que cette inclusion est valable sans restriction sur $\handwritten{\ell}$.

En effet, le membre de gauche est canoniquement isomorphe à $D(T_{\handwritten{\ell}}(A'_K)/T_{\handwritten{\ell}}(A'_K)^f)$, et est donc de ``rang'' (ou ``hauteur'') égal à $2n-2\handwritten{\alpha}-\handwritten{\mu}=\handwritten{\mu}$ (notations de (2.1.9), où par hypothèse $\handwritten{\lambda}=0$), qui est aussi celui de $T_{\handwritten{\ell}}(A_K)^t$. Comme l'inclusion

\[
T_{\handwritten{\ell}}(A_K)^t\ \handwritten{\subset}\ T_{\handwritten{\ell}}(A_K)\tag*{(*)}
\]

est \uline{stricte}, dans le sens qu'elle induit déjà un monomorphisme pour tous les termes du systèmes projectifs envisagés, il s'ensuit aisément que le premier membre de (*) est maximal parmi les sous-pro-groupes de Barsotti-Tate de $T_{\handwritten{\ell}}(A_K)$ ayant même rang que lui, de sorte qu'on a bien l'égalité de (5.2.7), qui équivaut à l'inclusion entre les termes extrêmes de 5.2.7 par le théorème d'orthogonalité.

5.3. Utilisant la correspondance par orthogonalité entre sous-pro-groupes de Barsotti-Tate stricts (au sens précisé plus haut) des pro-groupe\unsure{ś|s} duals $T_{\handwritten{\ell}}(A_K)$ et $T_{\handwritten{\ell}}(A'_K)$, on déduit formellement de (5.2.2) que l'accouplement (1.0.2) entre ces pro-groupes induit un accouplement

\[
T_{\handwritten{\ell}}(A_K)^{\mathrm{ab}}\times T_{\handwritten{\ell}}(A'_K)^{\mathrm{ab}}\handwritten{\longrightarrow} T_{\handwritten{\ell}}(\mathbb{G}_{mK})\tag*{(5.3.2)}
\]

% archive sheet number: 200
% archive numbering: 375 (pencil and printed label, foot).

\scanpage{202}

\[
X_0+\underbrace{\sum_{ij\geq1}f_{ij}X_0^{(i)}\sum_{0\leq k\leq j}T^{(j-k)}(-1)^k X_0^k}_{\displaystyle\sum_{\substack{i,j,k\\0\leq k\leq j\,;\ i,j\geq1}}f_{ij}T^{(j-k)}(-1)^k X_0^{i+k}}
\]

\[
\deleted{=\sum}\qquad
\left\{\begin{aligned}
\underline{i+k}&=\text{\textcircled{\ensuremath{\scriptstyle\unsure{\nu|v}}}} \geq 1\\
0&\leq\text{\textcircled{\ensuremath{\scriptstyle k}}}\leq\nu-1\\
j&\geq\mathrm{Sup}\,(k,1)
\end{aligned}\right.\quad
\begin{aligned}
i&=\nu-k\geq1\\
j&=k+\delta\qquad\delta=j-k\\
&\hspace{5em}\delta\geq0
\end{aligned}
\]

\[
\deleted{\sum_{\nu}X_0^\nu\sum_{\substack{0\leq k\leq\nu-1\\\deleted{j\geq\mathrm{Sup}(k,1)}\\\delta\geq0\\(k,\delta)\ne(0,0)}}(-1)^k f_{\unsure{i|j|\nu},k+\delta}T^{(\deleted{\unsure{j-k}}\delta)}}
\]

\[
\sum_{\nu\geq1}X_0^\nu
\underbrace{\sum_{\substack{0\leq k\leq\nu-1\\\delta\geq0\\\deleted{(k,\delta)\ne0}\\k+\delta\geq1}}(-1)^k f_{\nu-k,k+\delta}T^\delta}_{a_{1,\nu}}
\]

\[
\begin{aligned}
&\deleted{\unsure{\sum}}\\
A_0&=X_0\overbrace{\biggl(1+\sum_{\delta\geq1}f_{1,\delta}T^{(\delta)}\biggr)}^{a_{1,1}}+\sum\\
&\hspace{5em}(k=0,\nu=1)
\end{aligned}
\]

\noindent\rule{\linewidth}{0.3pt}

\[
\deleted{\varphi(\xi_1}\qquad\varphi\bigl(\deleted{\unsure{(}}X_0^{\unsure{0|o}}\bigr)
\]

% archive sheet number: 201

\scanpage{203}\mspage{- 56 -}\hfill IX

\uline{Lemme} 5.2.5. \uline{Soit} $u:G \handwritten{\longrightarrow} H$ \uline{un homomorphisme de pro}-$\handwritten{\ell}$-\uline{groupes de Barsotti-Tate sur} $S$. \uline{Si} $u_{\unsure{o|0}}:G_o \handwritten{\longrightarrow} H_o$ \uline{est nul}, \uline{alors} $u$ \uline{est nul}. \uline{Si} $u_o:G_o \handwritten{\longrightarrow} H_o$ \uline{est une isogénie} (resp. \uline{un isomorphisme}), $u$ \uline{est une isogénie} (resp. \uline{un isomorphisme}).

Nous appliquerons en effet la première assertion de 5.2.5 au morphisme composé

\[
T_{\handwritten{\ell}}(A^o)^t \handwritten{\longrightarrow} T_{\handwritten{\ell}}(A^o)^f \handwritten{\longrightarrow} D(T_{\handwritten{\ell}}(A^{\prime o})^f)\quad ,
\]

où le deuxième morphisme est celui déduit de l'accouplement (5.2.4.2), on trouve que ce composé est nul, donc que l'accouplement envisagé s'annule sur la partie torique à gauche. Par symétrie, il s'annule donc sur la partie torique à droite, donc il provient d'un accouplement sur les quotients, ou ``morceaux abéliens''

\[
T_{\handwritten{\ell}}(A^o)^{\mathrm{ab}}\times T_{\handwritten{\ell}}(A^{\prime o})^{\mathrm{ab}} \handwritten{\longrightarrow} T_{\handwritten{\ell}}(\mathbb{G}_m)\quad ,
\]

lesquels quotients sont encore des pro-groupes de Barsotti-Tate, comme on vérifie aisément. Il reste à prouver que le morphisme correspondant à cet accouplement

\[
T_{\handwritten{\ell}}(A^o)^{\mathrm{ab}} \handwritten{\longrightarrow} D(T_{\handwritten{\ell}}(A^{\prime o})^{\mathrm{ab}})
\]

est une isogénie, ce qui résulte \unsure{.} encore de 5.2.5.

Nous laisserons au lecteur la démonstration (facile) de 5.2.5, qui est un résultat valable sur toute base noethérienne connexe. Nous trouverons d'ailleurs plus bas (7.4.7), par une analyse un peu plus précise des biextensions $W_n$ des $(A^o_n,A^{\prime o}_n)$, une démonstration directe assez évidente des deux cas particuliers où nous avons eu à utiliser 5.2.5.

% archive sheet number: 202
% Additional foliation at the foot in reading orientation: 374 (twice).
